% P56 | 7 Conclusions | Conclusions: the American put
% guide: dev_notes/paragraph_guides/P56_conclusions_the_american_put.txt
The CRR tree is an explicit finite-difference scheme for the variational inequality, and the tree, a finite-difference solver and least-squares Monte Carlo agree on the price of the American put.
The perpetual put is the limit of the long-dated put, and the one case that seemed to contradict this, a thousand-year maturity on a coarse lattice, was an artefact of the lattice.
Early exercise adds a few percent at the money and more deep in the money, and it needs a positive interest rate for a put and a dividend for a call.
The boundary falls when volatility rises and rises when the interest rate rises, which is what the cost and the value of waiting predict.
The price does not depend on the real-world drift but the buyer's expected profit does, and a seller who does not hedge carries that dependence.
A wrong volatility corrupts the exercise rule, and the loss is asymmetric: exercising too early costs much more than exercising too late.
In practice the volatility input and the exercise rule should be monitored together, because an error in the first propagates into the second.
